\documentclass[a4paper,12pt]{article}

\usepackage[utf8]{inputenc}
\usepackage[T1]{fontenc}
\usepackage[spanish]{babel}
\usepackage{amsmath, amssymb}
\usepackage{graphicx}
\usepackage{float}
\usepackage{subcaption}
\usepackage{hyperref}
\usepackage{listings}
\usepackage{xcolor}

% Define colors for syntax highlighting
\definecolor{codegreen}{rgb}{0,0.6,0}
\definecolor{codegray}{rgb}{0.5,0.5,0.5}
\definecolor{codepurple}{rgb}{0.58,0,0.82}
\definecolor{backcolour}{rgb}{0.95,0.95,0.92}

\lstdefinestyle{pythonstyle}{
    language=Python,
    backgroundcolor=\color{backcolour},
    commentstyle=\color{codegreen},
    keywordstyle=\color{magenta},
    numberstyle=\tiny\color{codegray},
    stringstyle=\color{codepurple},
    basicstyle=\ttfamily\footnotesize,
    breakatwhitespace=false,
    breaklines=true,
    captionpos=b,
    keepspaces=true,
    numbers=left,
    numbersep=5pt,
    showspaces=false,
    showstringspaces=false,
    showtabs=false,
    tabsize=4
}
\lstset{style=pythonstyle}



\begin{document}

{\centering \textsf{TFG Pepe}\par}

\vspace{0.5in}

{\huge Luns 13 de Abril de 2026}

\section{Propuesta anteproxecto}

\subsection{Titulo}
holai lala $\pi$
\subsection{Breve descripción}
Este traballo xurde no berce dun momento de desenvolvemento teórico guiado pola experimentación e o empiricismo, planteándose como oposición. As Redes de Neuronas Artificiais que están a sufrir un esplendor en aplicacións carecen dun marco teórico para entendelas tanto os profesionais informáticos como o resto que traballen con elas. Por eso apostamos pola proposta de Yi Ma que a través duns principios claramente caracterizados (compresión a través de representacións que conserven a máxima información intra clase e a máxima discriminación entre case) derivánse unha familia de Redes Neuronais sobre as que se planetea este traballo.
\subsection{Obxectivos concretos}
Estudar como as \emph{cabezas de atención semánticas} poden apoiar a interpretabilidade no diagnóstico no dominio médico  \textbf{e} como pequenos cambios na rede, xurdidos da súa formulación, poden afectar ó rendemento da mesma.
\subsection{Método de traballo}
A metodoloxía é de tipo iterativo incremental.
\subsection{Fases principais}
\begin{enumerate}
	\item Lectura bibliográfica
	\item Acomodación do código para experimentos futuros e planteamento da tarefa inicial (predicción binaria de vaso).
	\item Planteamento de propostas e experimentación.
\end{enumerate}
\subsection{Material e medios necesarios}
\begin{itemize}
	\item GPU do varpa con \texttt{Python 3.8.10} e \texttt{rsync} instalados.
	\item Ordenador portatil

\end{itemize}


\section{Que é o master?}
\section{Como son os parches da rede preentrenada do Yi Ma?}


\begin{figure}[H]
\centering
\includegraphics[width=\linewidth]{figs/parches_grid_demo_01.pdf}
\end{figure}

\section{Experimentos con moitas épocas [en curso]}

Lancei os nove experimentos correspondentes ás combinacións de usar ou non as seguintes pérdidas:
\begin{itemize}
	\item $||{A}||_1 + ||b||_1$, a chamo L1
	\item $||x -A^{'}Ax||_{fro}$, a chamo reconstruccion
	\item $||AA^{'} - I||_{fro}$, a chamo ortogonal
\end{itemize}

Por tanto son 8 ($2^3$) experimentos a 200000 epócas.

De momento a que mellor foi (que para min foi mal) é: sen L1, con reconstruccion e con ortogonal. As figuras de abaixo teñen ganacia de 5.
\begin{figure}[H]
\centering
\includegraphics[width=\linewidth]{figs/buena_exp9_ganancia5.pdf}
\end{figure}

\begin{figure}[H]
\centering
\includegraphics[width=\linewidth]{figs/buena_exp9_cum.pdf}
\end{figure}


\end{document}
